\scanpage{138}
\mspage{83}
% archive sheet number: 137

\unsure{sort|dont} la richesse des intuitions liées à une telle notion. Peut-être, sans le savoir ai-je toujours traité $\mathrm{Cat}$ non comme une catégorie ``nue'', mais comme un site, ou plus \unsure{t\textsuperscript{t}|tt} (puisqu'il y a l'involution $X\to X^{\unsure{\circ|o}}$ qui ne peut être éludée), comme un ``bi-site'' -- on y a \uline{deux} manières de localiser, \deleted{\unsure{par}} soit par les $X_{/x}$ (top. de Zariski!) soit par les $\,{}_x\backslash X$ (top. co-Zariskienne). Il y a aussi la topologie lisse (toute proche sans doute de l'intuition de la \unsure{cat.} étale), et la topologie propre -- et on s'attend que la lisse donne ess. les $\unsure{\Rightarrow|=}$ \unsure{résultats} que la Zariskienne \deleted{\unsure{-- mais peut-être pas juste}} (vu que la topologie étale \deleted{\unsure{plus}} est égale -- il n'y a pas lieu d'en faire un plat\ldots) ; et la topologie propre les mêmes résultats (p.ex. la $=$ cohomologie ?) que la top. co-Zariskienne.

Il semble bien que dans \deleted{\unsure{toutes}} les catégories de modèles utilisées jusqu'à présent pour décrire $\mathrm{Hot}$ comme un localisé, on dispose de notions de localisation: $(\mathrm{Top})$, $\Delta^\wedge$ (ou $A^\wedge$, avec $A$ catégorie-test quelconque), $(\mathrm{Cat})$; et dans $\mathrm{Cat}$ on dispose même de \uline{deux} telles notions. Je pense par contre aux $\infty$-groupoïdes, c'est peut-être la seule catégorie de modèles à tel point proche déjà de $\mathrm{Hot}$ \unsure{elle-même}, qu'on n'y discerne pas de notion de localisation, et que l'idée même de \unsure{vouloir|modèles} \unsure{s'y localiser|s'y localise} (si ce n'est en ``localisant'' $\mathrm{Hot}$ lui-même) paraît quasiment
